\subsection*{The exact subgroup intersection}

\statement{B.28}{Proved}{The pure tetrahedral and cyclic intersection}
Put $P=P\Gamma_{0,12}$, choose the named $T=A_4$ and $D=C_5$
inside the geometric $A_5$, and put
\[
 H_T=C_P(T),\qquad H_5=C_P(D).
\]
In the full profinite completion of the pure mapping class group,
\[
 \mathcal K=C_{\widehat P}(A_5)
       =\widehat{H_T}\cap\widehat{H_5},\qquad
 H_T\simeq F_{13},\quad H_5\simeq F_6,
 \qquad H_T\cap H_5=1.
\]
Both displayed completions are their actual embedded closures in
$\widehat P$. The ordered parameter curves and their basepoint
choices are those of B.18. In particular the equality is an
intersection in one specified ambient group, not an abstract
intersection of two unrelated free profinite groups.
\begin{proof}
The group $A_4$ is maximal in $A_5$, and $D$ is not contained in it,
so $\langle T,D\rangle=A_5$. Centralizers of a generated subgroup
are the intersections of the two centralizers, in both categories.
B.30 proves the required full topology and density for $T$ by the
$V_4$ and $C_3$ prime-power steps. For $D$ the same argument needs
only the single $5$-subgroup case of Boggi--Zalesskii Theorem A:
the prime-power ramification-filling proof of symmetric CSP gives
full induced topology, and the prime-power normalizer theorem
followed by its open centralizer kernel gives density. Intersecting
with the ambient open pure group preserves both assertions.

The pure tetrahedral subgroup has index twelve in $F_T=F_2$,
and hence rank thirteen, as proved in B.30. For the cyclic group,
place its fixed poles at $0,\infty$ and write its two free orbits as
$\mu_5,t\mu_5$. Ordered continuation permits exactly
$t\in\mathbb P^1\setminus\{0,\infty,\mu_5\}$.
The covering normalizer sequence has quotient the point-pushing
group on the four-point quotient sphere. Pure continuation fixes
all labels, and its degree-five parameter cover is the stated
seven-punctured curve; its deck rotations are detected in the
puncture permutation quotient and introduce no pure finite kernel.
The discrete branched-cover normalizer map is faithful on this
pure subgroup. Thus its actual monodromy is $H_5\simeq F_6$.
This also proves the identification of its completion used above.
Finally B.7 computes the discrete centralizer of $A_5$ as $1$.
The intersection of the two embedded completions is precisely
the unresolved profinite centralizer, not the completion of
this last discrete intersection unless an additional theorem
is supplied.
\end{proof}
